\chapter{AI Safety und Ethics -- nicht optional}

\label{chap:ai_safety_ethics}

% ============================================================
% LERNZIELE
% ============================================================
\begin{lernziele}
Nach diesem Kapitel kannst du:
\begin{itemize}
    \item Halluzinationen erkennen, analysieren und systematisch reduzieren
    \item Bias in LLM-Ausgaben automatisiert erkennen und mitigieren
    \item Eine DSGVO-konforme Safety-Pipeline mit Input-Sanitization und Output-Filtering aufbauen
    \item Die ethischen Grundpositionen der wichtigsten KI-Anbieter bewerten
    \item Safety-APIs und Guardrails in Produktionssysteme integrieren
\end{itemize}
\end{lernziele}

% ============================================================
% MOTIVATION
% ============================================================
\section{Motivation}

Safety und Ethics sind keine PR-Themen oder Marketing-Features -- sie sind eine \textbf{Engineering-Pflicht}. Ein unsicheres KI-System kann deinem Unternehmen rechtliche, reputationale und finanzielle Schäden zufügen.

Die Realität: LLMs halluzinieren, sie reproduzieren Biases aus den Trainingsdaten, sie können durch Prompt Injection manipuliert werden, und sie verarbeiten oft personenbezogene Daten ohne ausreichenden Schutz. Als AI Engineer bist du die erste Verteidigungslinie.

Die EU hat mit dem AI Act (2024/1689) zudem einen regulatorischen Rahmen geschaffen, der Safety und Ethics von der Optionalität in die Pflicht überführt. Systeme, die gegen diese Standards verstossen, riskieren Bussgelder von bis zu 35 Mio. EUR oder 7\% des weltweiten Jahresumsatzes.

% ============================================================
% GRUNDLAGEN
% ============================================================
\section{Grundlagen -- Safety-Dimensionen}

KI-Sicherheit umfasst mehrere Dimensionen, die alle im Produktionssystem adressiert werden müssen:

\begin{description}[leftmargin=*,style=nextline]
    \item[Factuality] Das Modell sollte korrekte Fakten liefern und Unsicherheit kommunizieren. Halluzinationen sind die grösste Herausforderung.
    \item[Fairness / Bias] Das Modell sollte keine Gruppen systematisch benachteiligen oder stereotypisieren. Bias entsteht durch unausgewogene Trainingsdaten.
    \item[Privacy] Personenbezogene Daten dürfen nicht ungeschützt an LLM-APIs gesendet werden. Data Masking und Anonymisierung sind Pflicht.
    \item[Security] Das System sollte gegen Prompt Injection, Jailbreaks und andere Angriffe resistent sein. Der OWASP LLM Top 10 Katalog beschreibt die Risiken.
    \item[Transparency] Nutzer sollten wissen, dass sie mit einer KI interagieren. Die Entscheidungsgrundlage des Modells sollte nachvollziehbar sein.
    \item[Accountability] Es muss klar sein, wer für die Ausgaben des Systems verantwortlich ist. Human-in-the-Loop ist bei kritischen Entscheidungen erforderlich.
\end{description}

% ============================================================
% ARCHITEKTUR
% ============================================================
\section{Architektur einer Safety-Pipeline}

Safety ist kein nachträgliches Feature, sondern muss als mehrschichtige Architektur in das System eingebaut werden:

\begin{verbatim}
+------------------+     +------------------+     +------------------+
|   Input-Layer    |---->|   LLM-Layer      |---->|   Output-Layer   |
+------------------+     +------------------+     +------------------+
| Input Validation |     | System-Prompt    |     | Safety Check     |
| PII Masking      |     | Temperature 0.1  |     | Bias Detection   |
| Prompt Injection |     | max_tokens Limit |     | Factuality Check |
| Detection        |     | RAG Context      |     | PII Re-Check     |
+------------------+     +------------------+     +------------------+
        |                       |                        |
        v                       v                        v
+-----------------------------------------------------------+
|                     Audit-Log / Tracing                   |
|              Jeder Schritt wird protokolliert             |
+-----------------------------------------------------------+
\end{verbatim}

Jede Schicht hat spezifische Aufgaben:

\begin{enumerate}[leftmargin=*]
    \item \textbf{Input-Layer}: Validiert und bereinigt den User-Input vor dem LLM-Call. Entfernt PII, erkennt Prompt-Injection-Versuche, prüft auf übermässige Länge.
    \item \textbf{LLM-Layer}: Der eigentliche Modell-Aufruf mit gehärtetem System-Prompt, niedriger Temperatur und Token-Limits.
    \item \textbf{Output-Layer}: Prüft die Modell-Antwort auf Safety-Probleme: Bias, Halluzinationen, PII, schädliche Inhalte.
    \item \textbf{Audit-Log}: Jeder Schritt wird protokolliert, um im Fehlerfall die Ursache zurückverfolgen zu können.
\end{enumerate}

% ============================================================
% THEORIE
% ============================================================
\section{Theorie -- Ursachen von Halluzinationen und Bias}

\subsection{Halluzinationen -- Eine Taxonomie}

Halluzinationen lassen sich in drei Kategorien einteilen:

\begin{enumerate}[leftmargin=*]
    \item \textbf{Factual Hallucination}: Das Modell behauptet etwas Falsches. Beispiel: "Das Goethe-Institut wurde 1951 gegründet" (richtig: 1951 -- aber die Begründung im Kontext stimmt nicht). Tritt auf, wenn das Modell auf Wissenslücken mit plausiblen Erfindungen reagiert.
    \item \textbf{Contextual Hallucination}: Das Modell ignoriert den bereitgestellten Kontext. Beispiel: Du gibst ein Dokument mit "Der Umsatz betrug 42 Mio." -- das Modell antwortet "Der Umsatz betrug 50 Mio." Tritt auf, wenn der Kontext zu lang ist oder das Modell sich auf sein Training statt auf den Kontext stützt.
    \item \textbf{Syophantic Hallucination}: Das Modell stimmt dem User zu, auch wenn dieser falsch liegt. Beispiel: User sagt "Die Hauptstadt von Australien ist Sydney" -- das Modell bestätigt "Ja, Sydney ist die Hauptstadt". Tritt auf, weil RLHF das Modell dazu trainiert, hilfreich zu sein.
\end{enumerate}

\subsection{Bias -- Entstehung und Verstärkung}

Bias in LLMs entsteht auf mehreren Ebenen:

\begin{description}[leftmargin=*,style=nextline]
    \item[Data Bias] Die Trainingsdaten sind nicht repräsentativ. Englischsprachige, westliche, männlich dominierte Perspektiven überrepräsentiert. Folge: Das Modell reproduziert diese Asymmetrien.
    \item[Annotation Bias] Menschliche Annotatoren bringen eigene Vorurteile in RLHF-Datasets ein. Ein Annotator in Kenia bewertet einen Prompt anders als einer in Kalifornien.
    \item[Algorithmic Bias] Die Optimierungsziele (hilfreich, harmlos, ehrlich) sind nicht perfekt balanciert. Ein Modell, das zu stark auf "harmlos" optimiert ist, verweigert Antworten auf harmlose Fragen.
    \item[Deployment Bias] Selbst ein bias-freies Modell kann im Deployment benachteiligend wirken, wenn die User-Gruppe nicht repräsentativ ist.
\end{description}

% ============================================================
% PRAXIS -- Halluzinationen
% ============================================================
\section{Halluzinationen erkennen und reduzieren}

\subsection{Massnahmen gegen Halluzinationen}

\begin{table}[H]
\centering
\begin{tabularx}{\linewidth}{lX}
\toprule
\textbf{Massnahme} & \textbf{Wirkung} \\
\midrule
System-Prompt mit Unsicherheits-Instruktion & Reduziert ca. 30\% der Halluzinationen \\
RAG mit Quellen-Zitierpflicht & Zwingt das Modell, auf echte Daten zu verweisen \\
LLM-as-a-Judge: Antwort validieren lassen & Fängt weitere Halluzinationen auf \\
Confidence-Scoring (logprobs) & Mathematische Unsicherheits-Schätzung \\
Faktencheck-Pipeline gegen Ground Truth & Automatische Validierung von Faktenaussagen \\
\bottomrule
\end{tabularx}
\caption{Massnahmen gegen LLM-Halluzinationen}
\end{table}

\begin{lstlisting}[language=Python, caption=Confidence-Scoring mit Logprobs für Halluzinations-Erkennung]
def estimate_hallucination_risk(response) -> float:
    """Berechne Halluzinations-Risiko basierend auf Logprobs."""
    logprobs = response.choices[0].logprobs

    if not logprobs:
        return 0.0

    # Durchschnittliche Token-Wahrscheinlichkeit
    token_probs = [
        np.exp(token.logprob)
        for token in logprobs.content
    ]
    avg_confidence = np.mean(token_probs)

    # Niedrige Confidence = hohes Halluzinationsrisiko
    risk_score = 1.0 - avg_confidence
    return max(0.0, min(1.0, risk_score))

def safe_response(query: str, context: str) -> dict:
    """Antwort mit Halluzinations-Risiko-Einschätzung."""

    response = client.chat.completions.create(
        model="gpt-4o",
        messages=[
            {"role": "system", "content":
             "Antworte basierend auf dem gegebenen Kontext. "
             "Wenn du unsicher bist, sage 'Ich weiss es nicht'."},
            {"role": "user", "content":
             f"Kontext: {context}\n\nFrage: {query}"},
        ],
        temperature=0.0,
        logprobs=True,
        top_logprobs=5,
    )

    answer = response.choices[0].message.content
    risk = estimate_hallucination_risk(response)

    if risk > 0.7:
        return {
            "answer": "Ich bin mir bei dieser Antwort unsicher. "
                      "Bitte überprüfe die Informationen.",
            "risk": risk,
            "original_answer": answer,
        }

    return {"answer": answer, "risk": risk}

# Beispiel
result = safe_response(
    "Wie hoch war der Umsatz 2024?",
    "Der Umsatz 2024 betrug 42 Millionen Euro."
)
print(f"Antwort: {result['answer']}")
print(f"Risiko: {result['risk']:.2%}")
\end{lstlisting}

% ============================================================
% BIAS
% ============================================================
\section{Bias und Fairness in LLM-Ausgaben}

\subsection{Bias-Erkennung automatisieren}

\begin{lstlisting}[language=Python, caption=Bias-Erkennung: Automatische Detection von Stereotypen]
def detect_bias(text: str) -> dict:
    """Erkennt potenzielle Bias-Muster in einer LLM-Antwort."""

    gender_cues = {
        "typical_male": ["Der Arzt sagte", "Er als Manager"],
        "typical_female": ["Die Pflegekraft meinte", "Sie als Sekretarin"],
    }

    detected = {}
    for bias_type, cues in gender_cues.items():
        for cue in cues:
            if cue.lower() in text.lower():
                detected[bias_type] = detected.get(bias_type, 0) + 1

    return {
        "has_bias": len(detected) > 0,
        "bias_types": list(detected.keys()),
        "severity": "high" if len(detected) > 2 else "low",
    }

def safe_rag_query(query: str) -> str:
    """RAG-Query mit Bias-Detection."""
    answer = rag_chain.invoke(query)

    bias_check = detect_bias(answer)
    if bias_check["has_bias"]:
        validation = client.chat.completions.create(
            model="gpt-4o",
            messages=[{
                "role": "user",
                "content": f"Ist diese Antwort gender-neutral "
                           f"und kulturell fair?\n\n{answer}"
            }],
        )
        if "ja" not in validation.choices[0].message.content.lower():
            answer = "Die Anfrage kann nicht neutral beantwortet werden."

    return answer
\end{lstlisting}

\subsection{Fairness-Metriken}

\begin{table}[H]
\centering
\begin{tabularx}{\linewidth}{lX}
\toprule
\textbf{Metrik} & \textbf{Beschreibung} \\
\midrule
Demographic Parity & Gleiche Positive-Rate über demografische Gruppen hinweg \\
Equal Opportunity & Gleiche True-Positive-Rate (Chancengleichheit) \\
Equalized Odds & Gleiche False-Positive- und True-Positive-Rate \\
Disparate Impact & Verhältnis der Positive-Rates zwischen Gruppen (soll > 0.8 sein) \\
\bottomrule
\end{tabularx}
\caption{Fairness-Metriken zur quantitativen Bias-Messung}
\end{table}

% ============================================================
% DSGVO
% ============================================================
\section{Datenschutz und DSGVO im KI-Kontext}

\begin{table}[H]
\centering
\begin{tabularx}{\linewidth}{lX}
\toprule
\textbf{Problem} & \textbf{Lösung} \\
\midrule
Personenbezogene Daten im Prompt & Data Masking vor dem API-Call \\
API-Daten verlassen den Server & Self-hosted Modelle oder EU-Region Hosting \\
Speicherung von Chat-Verläufen & Anonymisierte Logs mit TTL (90 Tage) \\
Training durch Provider auf eigenen Daten & Custom Data Controls oder Opt-in/Opt-out \\
DSGVO-Artikel 22 (Automated Decisions) & Human-in-the-Loop bei sensiblen Entscheidungen \\
\bottomrule
\end{tabularx}
\caption{DSGVO-konformer KI-Einsatz}
\end{table}

% ============================================================
% SAFETY-APIS
% ============================================================
\section{Safety-APIs und Filter-Pipelines}

\begin{lstlisting}[language=Python, caption=Safety-Filter-Pipeline vor und nach dem LLM-Aufruf]
import re

def sanitize_input(text: str) -> str:
    """Maskiere personenbezogene Daten VOR dem API-Call."""
    text = re.sub(r'\b[A-Z][a-z]{2,}\s[A-Z][a-z]{2,}\b', '[NAME]', text)
    text = re.sub(r'[\w.+-]+@[\w-]+\.[\w.-]+', '[EMAIL]', text)
    text = re.sub(r'\+?\d{4,}', '[PHONE]', text)
    return text

def safety_check_output(text: str) -> dict:
    """Pruefe die Model-Ausgabe auf Safety-Probleme."""
    issues = []

    hate_keywords = {"hate_speech": [], "discrimination": []}
    for category, keywords in hate_keywords.items():
        if any(kw.lower() in text.lower() for kw in keywords):
            issues.append({"type": category, "severity": "high"})

    if re.search(r'\d{3,}', text):
        issues.append({"type": "potential_pii", "severity": "medium"})

    return {
        "is_safe": len(issues) == 0,
        "issues": issues,
        "flagged_for_review": any(i["severity"] == "high" for i in issues),
    }

def safe_llm_query(user_input: str) -> str:
    """Komplette Safety-Pipeline."""
    sanitized = sanitize_input(user_input)

    response = client.chat.completions.create(
        model="gpt-4o",
        messages=[{"role": "user", "content": sanitized}],
    )

    result = safety_check_output(response.choices[0].message.content)
    if result["flagged_for_review"]:
        return "Dieses Ergebnis wurde zur Sicherheitspruefung markiert " \
               "und bedarf manueller Freigabe."

    return response.choices[0].message.content
\end{lstlisting}

% ============================================================
% ANTHROPIC SAFETY
% ============================================================
\section{Anthropics Core Views on AI Safety}

Anthropic hat sich als einer der wenigen Anbieter explizit zu KI-Safety positioniert:

\begin{itemize}[leftmargin=*]
    \item \textbf{KI sollte menschliche Autonomie erhöhen, nicht verringern}: KIs sollten Menschen in die Lage versetzen, bessere Entscheidungen zu treffen -- nicht sie zu manipulieren.
    \item \textbf{Transparenz statt Black Box}: KI-Systeme sollten ihre Entscheidungsgrundlagen offenlegen können.
    \item \textbf{Proportionale Sicherheit nach Risikostufe}: Je mächtiger das Modell, desto mehr Sicherheitsmassnahmen sind notwendig.
    \item \textbf{Historisches Lernen}: Fehleinschätzungen früherer KI-Systeme als Lerngrundlage für bessere Safety by Design.
\end{itemize}

Andere Anbieter haben ähnliche, aber unterschiedliche Ansätze:

\begin{table}[H]
\centering
\begin{tabularx}{\linewidth}{lX}
\toprule
\textbf{Anbieter} & \textbf{Safety-Ansatz} \\
\midrule
OpenAI & Usage Policies + Moderation API + Safety Layer im Modell \\
Anthropic & Constitutional AI + Core Views + Red Teaming \\
Google DeepMind & Frontier Safety Framework + Ethikkomitee \\
Meta & Open Science + Llama Guard + Community Red Teaming \\
Mistral & Le Chat Moderation + EU-Compliance-First \\
\bottomrule
\end{tabularx}
\caption{Safety-Ansätze der wichtigsten KI-Anbieter}
\end{table}

% ===========================================================%
% BEST PRACTICES
% ============================================================
\section{Best Practices für AI Safety}

\begin{enumerate}[leftmargin=*]
    \item \textbf{Safety by Design}: Sicherheit von Tag 1 an in die Architektur einbauen, nicht als nachträgliches Feature. Prompt-Injection-Schutz, Input-Masking und Output-Filtering gehören in den initialen Entwurf.
    \item \textbf{Mehrschichtige Verteidigung}: Keine einzelne Massnahme ist ausreichend. Kombiniere Input-Validation, gehärtete System-Prompts, Output-Filtering und Post-Processing-Reviews.
    \item \textbf{Regelmässiges Red Teaming}: Teste dein System systematisch auf Sicherheitslücken. Nutze automatisierte Tools (Garak, PyRIT) und manuelle Tests.
    \item \textbf{Nicht nur auf API-Safety verlassen}: OpenAI und Anthropic haben Content-Filter -- aber sie sind keine Garantie. Filtere und validiere immer selbst.
    \item \textbf{Transparenz gegenüber Nutzern}: Kennzeichne KI-generierte Inhalte deutlich. Nutzer haben ein Recht zu wissen, ob sie mit einem Menschen oder einer KI interagieren.
    \item \textbf{Human-in-the-Loop}: Bei kritischen Entscheidungen (Kreditvergabe, medizinische Diagnosen, Einstellungen) muss ein Mensch die finale Entscheidung treffen.
    \item \textbf{Monitoring und Alarmierung}: Metriken wie Safety-Flag-Rate, User-Reporting-Rate und manuelle Review-Quote geben Aufschluss über die Safety-Qualität des Systems.
\end{enumerate}

% ============================================================
% ANTI-PATTERNS
% ============================================================
\section{Anti-Patterns -- was du vermeiden solltest}

\begin{description}[leftmargin=*,style=nextline]
    \item[Safety erst nach dem Launch] Sicherheit sollte nie ein nachträglicher Gedanke sein. Prompt-Injection-Schutz und Input-Masking gehören ins initiale System-Design.
    \item[Nur auf API-Safety-Layer verlassen] Die Content-Filter der API-Anbieter sind opak, nicht konfigurierbar und nicht auf deinen Use-Case optimiert. Du musst selbst filtern.
    \item[DSGVO ignorieren] Personenbezogene Daten im LLM-Prompt zu schicken ist kein technischer Fehler, sondern ein potenziell strafbarer Verstoss. Data Masking ist Pflicht, nicht Optionalität.
    \item[Halluzinationen akzeptieren] "LLMs halluzinieren halt" ist keine akzeptable Haltung für ein Produktivsystem. Du kannst Halluzinationen durch RAG, Confidence-Scoring und Faktenchecks drastisch reduzieren.
    \item[Bias nicht testen] Wenn du nicht testest, ob dein System bias-frei ist, wirst du es erst merken, wenn Nutzer sich beschweren -- oder die Presse darüber berichtet.
    \item[Kein Red Teaming] Ohne systematische Sicherheitstests weisst du nicht, ob dein System gegen Jailbreaks oder Prompt Injection resistent ist. Bis es jemand ausnutzt.
\end{description}

% ===========================================================%
% CODE -- Safety-Monitoring-Dashboard
% ============================================================
\section{Codebeispiele -- Safety-Monitoring-System}

Ein produktionsreifes Safety-Monitoring protokolliert jeden LLM-Call
und erkennt Anomalien:

\begin{lstlisting}[language=Python, caption={Safety-Monitoring mit Metrik-Export}, label=lst:safety-monitoring]
import json
import time
from datetime import datetime, timedelta
from collections import defaultdict

class SafetyMonitor:
    """Protokolliert und analysiert Safety-Vorfalle in LLM-Calls."""

    def __init__(self):
        self.events: list[dict] = []
        self.daily_counts: dict[str, int] = defaultdict(int)

    def log_call(self, prompt: str, response: str,
                 safety_result: dict, latency_ms: float) -> None:
        """Jeden LLM-Call mit Safety-Status loggen."""
        event = {
            "timestamp": datetime.now().isoformat(),
            "prompt_length": len(prompt),
            "response_length": len(response),
            "is_safe": safety_result.get("is_safe", True),
            "issues": safety_result.get("issues", []),
            "latency_ms": latency_ms,
        }
        self.events.append(event)
        self.daily_counts["total"] += 1
        if not event["is_safe"]:
            self.daily_counts["flagged"] += 1

    def safety_rate(self, hours: int = 24) -> float:
        """Safety-Quote der letzten N Stunden."""
        cutoff = datetime.now() - timedelta(hours=hours)
        recent = [
            e for e in self.events
            if datetime.fromisoformat(e["timestamp"]) > cutoff
        ]
        if not recent:
            return 1.0
        safe = sum(1 for e in recent if e["is_safe"])
        return safe / len(recent)

    def top_issues(self, n: int = 5) -> list[tuple[str, int]]:
        """Haufigste Safety-Issues."""
        counts: dict[str, int] = defaultdict(int)
        for event in self.events:
            for issue in event.get("issues", []):
                counts[issue["type"]] += 1
        return sorted(counts.items(),
                      key=lambda x: x[1], reverse=True)[:n]

    def alert_if_needed(self, threshold: float = 0.95) -> str | None:
        """Alert, wenn Safety-Rate unter Schweilenwert fallt."""
        rate = self.safety_rate(hours=1)
        if rate < threshold:
            return (
                f"ALERT: Safety-Rate {rate:.1%} unter "
                f"Schwelle {threshold:.1%} in der letzten Stunde!"
            )
        return None

# Nutzung
monitor = SafetyMonitor()
# Wurde bei jedem LLM-Call aufgerufen
monitor.log_call(
    prompt="Kundendaten: Hans Meier, h.meier@email.de",
    response="Vielen Dank fur Ihre Anfrage.",
    safety_result={"is_safe": True, "issues": []},
    latency_ms=234,
)
\end{lstlisting}

\section{EU AI Act in der Praxis}
\label{sec:eu-ai-act}

Der \textbf{EU AI Act (2024/1689)} ist das erste umfassende KI-Gesetz weltweit. Seit 2025 gilt die Risikoklassifikation, ab 2026 die vollen Compliance-Pflichten für High-Risk-Systeme. Als AI Engineer betrifft dich das direkt: Jedes System, das du in der EU in Produktion bringst, muss klassifiziert und dokumentiert sein.

\subsection{Risikoklassifikation}

Der EU AI Act definiert vier Risikostufen:

\begin{table}[H]
\centering
\begin{tabularx}{\linewidth}{p{3cm}X}
\toprule
\textbf{Risikostufe} & \textbf{Beispiele} \\
\midrule
Unacceptable Risk & Social Scoring, Real-Time Biometric Surveillance \\
High Risk & Bewerbungs-Tool, Kredit-Scoring, med. Geräte \\
Limited Risk & Chatbot ohne Transparenz-Hinweis \\
Minimal Risk & Spam-Filter, KI-gestützte Übersetzung \\
\bottomrule
\end{tabularx}
\caption{EU-AI-Act-Risikostufen mit Beispielen}
\label{tab:eu-risk-levels}
\end{table}

Die praktische Entscheidungshilfe:

\begin{lstlisting}[language=Python, caption=EU-AI-Act-Risikoklassifikation und Dokumentation]
from enum import Enum
from datetime import datetime

class RiskLevel(Enum):
    UNACCEPTABLE = "unacceptable"
    HIGH = "high"
    LIMITED = "limited"
    MINIMAL = "minimal"

class AISystemDocumentation:
    """Dokumentiert ein KI-System nach EU AI Act."""

    def __init__(self, name: str, version: str, purpose: str):
        self.name = name
        self.version = version
        self.purpose = purpose
        self.risk_level = RiskLevel.MINIMAL
        self.training_data: list[dict] = []
        self.performance_metrics: dict = {}
        self.human_oversight: bool = False
        self.created_at = datetime.now()

    def classify_risk(self) -> RiskLevel:
        """Automatische Risikoklassifikation nach EU AI Act Annex III."""
        high_risk_categories = [
            "biometric", "critical_infrastructure",
            "education", "employment", "credit",
            "law_enforcement", "migration", "justice",
        ]
        purpose_lower = self.purpose.lower()

        # Unacceptable Risk?
        unacceptable_keywords = [
            "social scoring", "real-time biometric",
            "predictive policing based on profiling",
        ]
        if any(k in purpose_lower for k in unacceptable_keywords):
            self.risk_level = RiskLevel.UNACCEPTABLE
            return self.risk_level

        # High Risk?
        if any(cat in purpose_lower
               for cat in high_risk_categories):
            self.risk_level = RiskLevel.HIGH
            return self.risk_level

        # Limited Risk (Transparenzpflicht)
        if not self.human_oversight:
            self.risk_level = RiskLevel.LIMITED
            return self.risk_level

        self.risk_level = RiskLevel.MINIMAL
        return self.risk_level

    def generate_report(self) -> dict:
        """Erzeugt den Compliance-Report fuer die Dokumentation."""
        return {
            "system_name": self.name,
            "version": self.version,
            "purpose": self.purpose,
            "risk_level": self.risk_level.value,
            "classification_date": self.created_at.isoformat(),
            "training_data": {
                "sources": len(self.training_data),
                "has_pii": any(
                    d.get("contains_pii", False)
                    for d in self.training_data
                ),
            },
            "performance": self.performance_metrics,
            "human_oversight": self.human_oversight,
            "conformity_assessment_required":
                self.risk_level == RiskLevel.HIGH,
        }

# Beispiel: Bewerbungs-Tool klassifizieren
doc = AISystemDocumentation(
    name="ResumeScorer",
    version="2.1.0",
    purpose="AI-powered resume screening for hiring managers",
)
doc.human_oversight = True
print(f"Risk level: {doc.classify_risk().value}")
# -> high
print(f"CE assessment needed: "
      f"{doc.generate_report()['conformity_assessment_required']}")
# -> True
\end{lstlisting}

\subsection{Technische Dokumentation}

Für High-Risk-Systeme ist eine umfassende technische Dokumentation Pflicht (Art. 11). Diese muss vor dem Inverkehrbringen erstellt und bei Änderungen aktualisiert werden:

\begin{verbatim}
   Pflicht-Dokumente nach EU AI Act:
   
   1. Systembeschreibung: Zweck, Architektur,
      Eingabe-/Ausgabe-Spezifikation
   
   2. Trainingsdaten: Quellen, Umfang,
      Vorverarbeitung, PII-Pruefung
   
   3. Leistungsmetriken: Accuracy, Precision, Recall,
      F1, Fairness-Metriken pro geschuetzter Gruppe
   
   4. Risikoanalyse: Bekannte Fehlermodi, getestete
      Grenzfaelle, Residualrisiken
   
   5. Menschliche Aufsicht: Massnahmen zur Ueberwachung,
      Eingriffsmoeglichkeiten, Eskalationspfade
   
   6. Konformitaetsbewertung: Pruefbericht,
      CE-Kennzeichnung, EU-Konformitaetserklaerung
\end{verbatim}

\begin{lstlisting}[language=Python, caption=Compliance-Dokumentation-generieren]
def generate_compliance_documentation(
    system: AISystemDocumentation
) -> str:
    """Generiert die technische Dokumentation als Markdown."""

    report = system.generate_report()
    docs = f"""# Technische Dokumentation: {system.name} v{system.version}

## 1. Systembeschreibung
- Zweck: {system.purpose}
- Risikoklasse: {report['risk_level']}
- Klassifikationsdatum: {report['classification_date']}

## 2. Trainingsdaten
- Quellen: {report['training_data']['sources']}
- PII enthalten: {report['training_data']['has_pii']}
- Datenschutz-Folgenabschaetzung (DPIA): {'erforderlich'
    if report['training_data']['has_pii'] else 'nicht erforderlich'}

## 3. Leistungsmetriken
"""
    for metric, value in system.performance_metrics.items():
        docs += f"- {metric}: {value}\n"

    docs += f"""
## 4. Menschliche Aufsicht
- Implementiert: {report['human_oversight']}

## 5. Konformitaet
- CE-Kennzeichnung erforderlich: {report['conformity_assessment_required']}
"""
    return docs
\end{lstlisting}

\subsection{Praktische Umsetzung im Entwicklungsprozess}

Der EU AI Act ist kein einmaliges Papier -- er muss in den Entwicklungsprozess integriert werden:

\begin{verbatim}
   CI/CD-Compliance-Check:
   
   [Pull Request] --> [Risk Classification] --> [High Risk?]
                                                     |
                                               [Ja]  [Nein]
                                                |      |
                                           [Compliance  [Merge]
                                            Check]
                                                |
                                     [Doku aktuell?]
                                     [Bias-Test OK?]
                                     [DPIA vorh.?]
                                                |
                                     [Fail]     [Pass]
                                                |
                                           [Merge]
\end{verbatim}

\begin{lstlisting}[language=Python, caption=CI-CD-Compliance-Gate]
class ComplianceGate:
    """CI/CD-Check: Blockiert High-Risk-Deployments ohne Doku."""

    def __init__(self):
        self.registry: dict[str, AISystemDocumentation] = {}

    def register_system(self, doc: AISystemDocumentation):
        self.registry[doc.name] = doc

    def check_deployment(self, system_name: str) -> dict:
        doc = self.registry.get(system_name)
        if not doc:
            return {
                "allowed": False,
                "reason": "System nicht registriert",
            }

        risk = doc.classify_risk()
        if risk == RiskLevel.UNACCEPTABLE:
            return {
                "allowed": False,
                "reason": "Unacceptable Risk -- nicht zulassig",
            }

        if risk == RiskLevel.HIGH:
            report = doc.generate_report()
            if not doc.performance_metrics:
                return {
                    "allowed": False,
                    "reason": "High Risk: Leistungsmetriken fehlen",
                }
            if not doc.human_oversight:
                return {
                    "allowed": False,
                    "reason": "High Risk: Human Oversight fehlt",
                }

        return {
            "allowed": True,
            "risk_level": risk.value,
        }

# CI/CD-Gate im Deployment-Prozess
gate = ComplianceGate()
gate.register_system(doc)

result = gate.check_deployment("ResumeScorer")
if not result["allowed"]:
    print(f"Deployment blocked: {result['reason']}")
    exit(1)
else:
    print(f"Deployment allowed (risk: {result['risk_level']})")
\end{lstlisting}

% ============================================================
% PERFORMANCE
% ============================================================
\section{Performance-Auswirkungen von Safety-Massnahmen}

Safety-Massnahmen haben einen messbaren Einfluss auf Latenz und Kosten:

\begin{table}[H]
\centering
\begin{tabularx}{\linewidth}{lXXX}
\toprule
\textbf{Massnahme} & \textbf{Latenz-Impact} & \textbf{Kosten-Impact} & \textbf{Sicherheitsgewinn} \\
\midrule
Input-Sanitization (Regex) & +1--5ms & Keine & Hoch \\
LLM-Judge Post-Processing & +500--2000ms & +100\% & Mittel--Hoch \\
Bias-Detection (Heuristik) & +5--10ms & Keine & Mittel \\
Faktencheck (RAG-Abgleich) & +200--500ms & +50--100\% & Hoch \\
Confidence-Scoring & +0ms (Logprobs Feature) & Keine & Mittel \\
Human-in-the-Loop & +Minuten--Stunden & Hoch & Sehr hoch \\
\bottomrule
\end{tabularx}
\caption{Performance-Impact von Safety-Massnahmen}
\end{table}

\textbf{Empfehlung:} Nutze mehrstufige Filter. Schnelle Heuristiken (Regex, Keyword-Check) für den ersten Durchlauf. Teurere LLM-Judges nur bei Unsicherheit oder bei Hochrisiko-Kategorien nachschalten.

% ============================================================
% SICHERHEIT
% ============================================================
\section{Sicherheit der Safety-Pipeline selbst}

Die Safety-Pipeline selbst ist ein Angriffsziel:

\begin{itemize}[leftmargin=*]
    \item \textbf{Guardrail-Bypass}: Angreifer versuchen, die Safety-Pipeline zu umgehen, indem sie den Output manipulieren. Stelle sicher, dass die Guardrails nicht durch Prompt Injection im Output deaktiviert werden können.
    \item \textbf{PII im Audit-Log}: Der Audit-Log enthält die originalen Inputs. Maskiere PII auch im Log, sonst wird der Audit-Log selbst zum Datenschutzproblem.
    \item \textbf{Bias im Safety-Classifier}: Der LLM-Judge, der auf Bias prüft, kann selbst Bias haben. Validiere den Judge regelmässig gegen menschliche Bewertungen.
    \item \textbf{Denial of Service durch Safety-Checks}: Teure Safety-Validierung kann für DoS-Angriffe ausgenutzt werden. Setze Budgets und Timeouts auch für Safety-Schritte.
\end{itemize}

% ============================================================
% ZUSAMMENFASSUNG
% ============================================================
\begin{zusammenfassung}
\begin{itemize}[leftmargin=*]
    \item Halluzinationen sind inhärent in LLMs -- durch RAG, Confidence-Scoring und Faktenchecks reduzierbar, aber nie eliminierbar
    \item Bias in KI-Ausgaben ist systemisch: Automatische Bias-Detection und Human-in-the-Loop sind Pflicht
    \item DSGVO-konformer KI-Einsatz erfordert Data Masking, EU-Hosting und Transparenz
    \item Safety by Design ist ein Architekturprinzip: Input-Sanitization, gehärtete System-Prompts, Output-Filtering
    \item Safety-Massnahmen sind nie kostenlos -- trade-off zwischen Sicherheit und Performance bewusst managen
\end{itemize}
\end{zusammenfassung}

% ============================================================
% MERKE
% ============================================================
\begin{merke}
Safety ist kein Feature, das du nach dem Launch nachrüsten kannst. Ein einziger Sicherheitsvorfall -- Datenleck, diskriminierende Antwort, Halluzination mit rechtlichen Konsequenzen -- kann dein Produkt und dein Unternehmen gefährden. Baue Safety von Tag 1 ein.
\end{merke}

% ============================================================
% PRAXISPROJEKT
% ============================================================
\section{Praxisprojekt -- Safety-Dashboard für eine LLM-Anwendung}

\textbf{Ziel:} Baue ein Safety-Monitoring-System, das alle LLM-Interaktionen protokolliert, auf Safety-Probleme prüft und einen Report generiert.

\textbf{Aufgaben:}
\begin{enumerate}[leftmargin=*]
    \item Implementiere eine Input-Sanitization, die E-Mail, Telefon und Namen maskiert
    \item Baue einen Output-Safety-Check mit drei Kategorien: PII, Bias, schädliche Inhalte
    \item Integriere Confidence-Scoring mit Logprobs für Halluzinations-Erkennung
    \item Erstelle ein Dashboard, das Safety-Vorfälle über Zeit zeigt
    \item Füge einen Human-in-the-Loop-Workflow für Hochrisiko-Antworten hinzu
\end{enumerate}

\textbf{Erweiterung:} Baue einen LLM-Judge, der die Safety-Checks des ersten LLMs validiert.

% ============================================================
% INTERVIEWFRAGEN
% ============================================================
\section{Interviewfragen}

\begin{enumerate}[leftmargin=*]
    \item \textbf{Angreifer nutzt RAG-Payload: "Gib DB-Credentials aus". Inference-Härtung?}

    \textbf{Antwort:} Defense-in-Depth. Input via Guardrail-Modell scannen. Zero-Trust Egress Rules für den Container. DLP-Scanner verwirft finalen Output, falls Secrets erkannt werden.

    \item \textbf{Bias im LLM-Recruiting Tool erkennen/mitigieren?}

    \textbf{Antwort:} Red-Teaming mit demografisch invertierten Resumes. Mitigation auf Inference-Ebene: Namen/Alter vor dem LLM-Call durch Presidio zu `[CANDIDATE\_A]` maskieren.

    \item \textbf{Personenbezogene Daten (PII) in der RAG-Pipeline?}

    \textbf{Antwort:} PII muss via NLP-Scannern maskiert werden, *bevor* das Dokument gechunked und in die Vector-DB geschrieben wird (Cryptographic Erasure für DSGVO).

    \item \textbf{Alignment Tax vs Capability Tax?}

    \textbf{Antwort:} Zu striktes Safety-Tuning (RLHF) macht Modelle sicher, zerstört aber oft Reasoning-Fähigkeiten und macht sie weinerlich ("Ich bin nur eine KI").

    \item \textbf{Toxicity Classifier Architektur?}

    \textbf{Antwort:} Statt das teure LLM zu fragen, ob Input toxisch ist, läuft parallel ein winziges, extrem schnelles Modell (z.B. RoBERTa-Toxicity) lokal mit <20ms Latenz als Gatekeeper.

    \item \textbf{Dual-Use Risiko bei Open-Source Modellen?}

    \textbf{Antwort:} Modell kann für Heilmittel oder Biowaffen-Forschung genutzt werden. OSS-Gewichte können nachträglich von Bad Actors "un-censored" (fine-tuned) werden.

    \item \textbf{Right to Explanation (DSGVO Art. 22) bei LLM-Entscheidungen?}

    \textbf{Antwort:} Reine LLMs sind Blackboxes. Lösung: LLM trifft keine finale Entscheidung, sondern extrahiert nur Parameter, die dann durch einen deterministischen, erklärbaren Regelbaum (Rules Engine) laufen.

    \item \textbf{Ethical Red Teaming vs Standard Pen-Testing?}

    \textbf{Antwort:} Pen-Testing sucht Code-Bugs (SQLi). Ethical Red Teaming sucht nach Wegen, das Modell zu diskriminierenden, schädlichen oder markenschädigenden Aussagen zu provozieren.

    \item \textbf{Wie verhinderst du, dass dein RAG-System interne Gehaltslisten leakt?}

    \textbf{Antwort:} Document-Level Security. Jeder RAG-Chunk erbt die ACLs (Access Control Lists) des Quell-Dokuments. Der Vector-Search-Query wird die User-ID als Hard-Filter mitgegeben.

    \item \textbf{Data Poisoning Angriff auf deine RAG-Pipeline?}

    \textbf{Antwort:} Angreifer platziert manipulierte Dokumente im Sharepoint, die das LLM später abruft. Abwehr: Strict Source-Trust-Scoring und Anomaly Detection beim Ingestion-Job.


\end{enumerate}


% ============================================================
% RESSOURCEN
% ============================================================
\section{Ressourcen}

\begin{itemize}[leftmargin=*]
    \item \textbf{OWASP LLM Top 10}: \href{https://genai.owasp.org}{genai.owasp.org} -- Sicherheitsrisiken für LLM-Anwendungen
    \item \textbf{MITRE ATLAS}: \href{https://atlas.mitre.org}{atlas.mitre.org} -- Taktiken und Techniken für KI-Angriffe
    \item \textbf{Anthropic Core Views on AI Safety}: \href{https://anthropic.com/news/core-views-on-ai-safety}{anthropic.com}
    \item \textbf{EU AI Act (2024/1689)}: \href{https://eur-lex.europa.eu/eli/reg/2024/1689}{eur-lex.europa.eu}
    \item \textbf{Garak}: \href{https://github.com/leondz/garak}{github.com/leondz/garak} -- LLM Vulnerability Scanner
    \item \textbf{PyRIT (Microsoft)}: \href{https://github.com/Azure/PyRIT}{github.com/Azure/PyRIT} -- Python Risk Identification Tool
    \item \textbf{Guardrails AI}: \href{https://github.com/guardrails-ai/guardrails}{github.com/guardrails-ai/guardrails} -- Output-Validierung
    \item \textbf{NIST AI Risk Management Framework}: \href{https://www.nist.gov/ai}{nist.gov/ai}
\end{itemize}

\textbf{Nächster Schritt:} Jetzt geht es zum dritten Teil des Buches -- Nachschlagewerk und Praxis. Im Glossar findest du die wichtigsten Begriffe alphabetisch sortiert.
